% Cote 119: offprints, reports, courses, CV and letters (typed and printed part): archive cote 119, site page 182 (PDF page 183).
% typescript. Draft compared: none.
% AI-transcribed by Claude Opus 5.5 (claude-opus-5-5, adaptive thinking, effort high), Message Batches API, batch r119_full (msgbatch_01DwZNm5Y7YsKBWyAN4sFs8w), 2026-11-13; not yet checked by eye.
\scanpage{119/182}
UNIVERSITE DES SCIENCES \& TECHNIQUES \qquad ANNEE UNIVERSITAIRE 1973-74

DU LANGUEDOC

Mathématiques

M. P. II

FONCTIONS HOLOMORPHES

\uline{(Suite et fin)}

6.10. \uline{Corollaire}

Si $f(z)$ est holomorphe dans $U$, on a $\int_{\gamma} f(z)\,dz = 0$ pour tout che-%
min fermé $\gamma$ de $U$ qui est homotope à un point dans $U$.

\uline{Démonstration} : la forme $f(x)\,dz$ est fermée (théorème de Cauchy 4.3). On ap-%
plique 6.8.

On peut donner aussi une forme plus générale de la proposition 5.4.

6.1.1. \uline{Proposition}.

Soit \unsure{$f$} une fonction holomorphe dans un ouvert $U$. Soit $a \in U$ et
soit $\gamma$ un chemin fermé de $U$ ne passant pas par $a$ et homotope à un point
dans $U$. On a alors
\[ (6.11.1) \qquad \frac{1}{2\,\pi\, i}\int_{\gamma} \frac{f(z)\,dz}{z-a} = I(\gamma,a)\ \ f(a) \]
où $I(\gamma,a)$ est défini en (5.1).

\uline{Démonstration} :la démonstration est identique à celle de (5.4) $g(z)$ étant dé-%
fini par $\dfrac{f(z) - f(a)}{z-a}$ si $z \neq a$ et par $g(a) = f'(a)$ le théorème 4.4 montre
que $g(z)\,dz$ est fermée et le théorème 6.8. montre que $\int_{\gamma} g(z)\,dz = 0$ d'où
le résultat par définition de l'indice.

7. \uline{Principe du maximum}

(7.1.) \uline{Inégalités de Cauchy}

Soit $f$ une fonction holomorphe dans la boule ouverte $B(z_o,\rho)$ et
$r$ un nombre tel que $0 < r < \rho$. On a vu dans (5.5) que pour $|z - z_o| < \rho$ on
a:

\mspage{\hand{182}}
